% TOP_PROBLEM: {"id": 4600044, "title": "Growth of jointly periodic points II", "category_id": 11, "category": {"id": 11, "name": "dynamical_systems", "display_name": "Dynamical Systems", "description": "Problems about long-term behavior of deterministic systems, Hamiltonian dynamics, and periodic orbits.", "slug": "dynamical-systems", "order_index": 11, "created_at": "2026-07-31T15:26:25.671Z"}, "status": "open", "problem_number": "AMR-045-0044", "set_id": 13}
% FIRST_POSTED: 2026-10-01
% ENTRY_KIND: statement_only
% UnsolvedMath ID 4600044; source catalogue code AMR-045-0044
% !TeX program = lualatex
% Definition and sources only; this file is not an LLM solution attempt.
\documentclass[11pt,a4paper]{article}
\usepackage[margin=25mm]{geometry}
\usepackage{fontspec}
\setmainfont{lmroman10-regular.otf}[BoldFont=lmroman10-bold.otf,
 ItalicFont=lmroman10-italic.otf,BoldItalicFont=lmroman10-bolditalic.otf]
\usepackage{amsmath,amssymb,mathtools}
\usepackage{unicode-math}
\setmathfont{latinmodern-math.otf}
\AtBeginDocument{\let\setminus\smallsetminus}
\usepackage{xurl,hyperref}
\hypersetup{colorlinks=true,urlcolor=blue}
\setcounter{secnumdepth}{0}
\setlength{\parindent}{0pt}
\setlength{\parskip}{5pt}
\setlength{\emergencystretch}{3em}
\begin{document}
\section*{Growth of jointly periodic points II}\label{4600044}
{\small UnsolvedMath ID 4600044 · AMR-045-0044 · Dynamical Systems · AMR Open Problem Lists}
\subsection{Definitions and mathematical statement}
Let $N\ge2$ and $X_N=\{0,\ldots,N-1\}^{\mathbb Z}$ with
left shift $\sigma_N$. Consider any surjective continuous map
$F:X_N\to X_N$ commuting with the shift. Such an $F$ is a
one-dimensional cellular automaton and is determined by a
finite-window local rule. Define
\[
 J_k(F)=\{x:\sigma_N^k x=x,\quad
                   F^m x=x\text{ for some integer }m\ge1\},
 \qquad \nu(F)=\limsup_{k\to\infty}|J_k(F)|^{1/k}.
\]
The question is whether the universal lower bound
\[
                           \nu(F)\ge\sqrt N
\]
holds for every alphabet size and every such $F$.

The count uses configurations with spatial period dividing $k$,
with arbitrary positive temporal period. Distinct configurations
are counted separately even when they belong to the same shift
orbit or the same temporal cycle. Thus the ambient count at
spatial period $k$ is $N^k$. The limsup allows oscillation as
$k$ varies; the claim does not require a bound
$|J_k(F)|\ge N^{k/2}$ for every $k$.

Equivalently, for each $\varepsilon>0$ there must be infinitely
many $k$ with $|J_k(F)|^{1/k}>\sqrt N-\varepsilon$.
Surjectivity is assumed only for the full configuration space.
The proposed square-root threshold is independent of the local
rule's radius, and is stronger than asking merely that
$\nu(F)>1$ for each individual automaton.
\subsection{Short English statement}
Is the exponential growth rate of jointly periodic configurations always at least the square root of the alphabet size?
\subsection{Sources}
\begin{itemize}
\item[S1] ULAM.AI and the UnsolvedMath Contributors, supplied problems.json, record 4600044 (AMR-045-0044), AMR Open Problem Lists. Catalogue identity and source transcription; CC BY 4.0.
\url{https://huggingface.co/datasets/ulamai/UnsolvedMath}.
\item[S2] Boyle - Open Problems in Symbolic Dynamics (2008), Problem/Question/Conjecture 25.4. Original source indexed by AMR-045-0044.
\url{https://www.math.umd.edu/~mboyle/papers/openfinalsub3nov2007.pdf}.
\end{itemize}
\end{document}
